%!TEX program = lualatex
\documentclass[ngerman,headheight=70pt]{scrartcl}
\usepackage[ngerman]{babel}
\usepackage{blindtext}
\usepackage[utf8]{luainputenc}
\usepackage{scrlayer-scrpage}
\usepackage[T1]{fontenc}
\usepackage{fontspec}
\usepackage{multicol}
\usepackage{enumitem}
\usepackage[german=quotes]{csquotes}
\usepackage[bottom=2cm,footskip=8mm]{geometry}
\setmainfont{Source Sans Pro}
\MakeOuterQuote{"}

\setlength{\parindent}{0pt}
\setlength{\parskip}{6pt}

\pagestyle{plain.scrheadings}
\thispagestyle{scrheadings}
\lohead{\vspace*{-1.5cm}\includegraphics[scale=0.4]{../images/stupa_logo_rot.pdf}}
\rohead{\\Präsidium des\\Studierendenparlaments}

\lofoot[Universität Hamburg · Präsidium des Studierendenparlaments\\\hrule Von-Melle-Park 5 · D-20146 Hamburg]{Universität Hamburg · Präsidium des Studierendenparlaments\\\hrule Von-Melle-Park 5 · D-20146 Hamburg}
\cofoot[]{}
\cohead[\thepage]{}

\renewcommand{\thesubsection}{TOP \arabic{subsection}}
\renewcommand{\thesubsubsection}{\arabic{subsubsection}.}
\setcounter{subsection}{-1}

\begin{document}
    UHH · StuPa-Präsidium · Von-Melle-Park 5 · D-20146 Hamburg

    %%% PAGE ONE (EDIT LATER) %%%%
    % Hier den Titel des Protokolls eintragen
    \section*{Protokoll der 14. Sitzung des Studierendenparlaments vom 26. Januar 2017}

    % Protokollanten eintragen
    \textbf{Protokoll: Gunhild Berdal}\\
    \textbf{Ort: VMP5 0077}\\
    \textbf{Beginn: 18.35 Uhr}\\
    \textbf{Ende: 22.50 Uhr}

    % Wohl die komplizierteste Stelle im Protokoll
    % Als Faustregel sollten nicht mehr als 3 Namen auf einer Zeile stehen
    % Einfach an dieser Vorlage orientieren
    % Im Idealfall ist hier nicht viel zu ändern, da kaum jemand fehlt.
    % Rücktritte müssen natürlich pro Protokoll angepasst werden
    \vspace{0.5cm}
    \begin{tabular}{ll}
        Anwesend: & \\
            RCDS (5 Sitze): & Ramon Weilinger \\
             CampusGrün (14 Sitze): & Laura Franzen, Geoffrey Youett, Elena Rysikova, \\
                                   & Philipp Droll, Yasemin Günther, Melf Johannsen, \\
                                   & Sonja Lattwesen, Mario Moldenhauer, Christian Lagod, \\
                                   & Jim Martens, Svenja Horn, Damian Hofmann, \\
                                   & Martin Sievert  \\
             BGZM-Liste (2 Sitze) : & \\
             WiWi (2 Sitze): & \\
             Unicorns (5 Sitze): &  \\
             Liste LINKS (3 Sitze): & Gunhild Berdal, Till Petersen, Sinah Mielich \\
             HWP (2 Sitze): &  \\
             MIN (4 Sitze): &  \\
             SDS* (3 Sitze): & Mena Winkler, Jacob Petersein, Artur Brückmann \\
             Bart-LISTE (2 Sitze): & Timo Zeimet, Malte Peters \\
             LHG (1 Sitz): & Tobias Heisig \\
             harte zeiten (1 Sitz): & Tobias Berking \\
             Jura (1 Sitz): & \\
             AL (2 Sitze): & Henri Weber \\
            & \\
        Entschuldigt: & Jennifer Maack (RCDS), Armin Günther (CG), Jakob Pape, \\
                      & Karen Martirosian (BGZM), Johannes Peplow, Andreas Hartkamp (Unicorns), \\
                      & Jan Detampel (MIN), Karima Schulze (AL)\\
                      &\\
        Unentschuldigt abwesend: & Kolja Kolb, Ramin Shakiba, Benjamin Welling (RCDS), \\
                                & Claas-Friso Hente, Charleen Lorenz (WiWi), Katharina Kucza, \\
                                & Annkathrin Löffler, Marielle Hermstrüwer (Unicorns), Hakan Tasdemir, \\
                                & Gökce Celik (HWP), Ailina Salten, Lotte Rullkötter, \\
                                & Nabil Jabareen (MIN), Maurice Martens (Jura)\\
        Rücktritte: & \\
        Mandatsverluste: & \\
    \end{tabular}
    \newpage
    %%% Vorgeschlagene Tagesordnung %%%
    \underline{Vorgeschlagene Tagesordnung}
    \begin{enumerate}[label={\textbf{Top \theenumi}},leftmargin=*]
        \item RBCS-Bestätigung (10 Min.)
        \begin{enumerate}
            \item Bestätigung der Wahlniederschrift
            \item Bestätigung der Referent*innen
        \end{enumerate}
        \item VG-Wort (V1617-047) (30 Min.)
        \item Englischsprachige Wahlwerbung (V1617-048) (30 Min.)
        \item Wahlauswertung (60 Min.)
        \item Gutachten des Wissenschaftsrats (V1617-050) (30 Min.)
        \item Verschiedenes (5 Min.)
    \end{enumerate}

    \newpage

    % Hier beginnt Teil A. Bei Fortsetzungen kommt hier nur der Teil über
    % die Feststellung der endgültigen Fassung der Tagesordnung.
    % Für sonstige Sitzungen ist dies bereits das Gerüst für die Formalia.
    % Man muss nur noch die protokollierten Inhalte eintragen.

    {\Large \underline{Teil A}}

    \subsection{Formalia}

    \subsubsection{Geschäftsbericht Präsidium}

    Um 18.35 Uhr eröffnet Ramon die Sitzung und berichtet: Die StuPa-Wahlen
    seien erfolgreich verlaufen. 6550 Studierende haben sich an der Wahl
    beteiligt, das entspricht 15,3\% der 42000 Wahlberechtigten. Ungültige
    Stimmen waren bislang nicht zur Wahlbeteiligung gezählt, jetzt sind sie
    mitgezählt. Dennoch seien 10\% ungültige Stimmen eine immens hohe Zahl.

    Das Präsidium har mehrheitlich beschlossen, zur konstituierenden Sitzung am
    6. April einzuladen.

    \subsubsection{Anfragen an das Präsidium}

    Keine.

    \subsubsection{Geschäftsbericht AStA}

    Der AStA sieht die "Empfehlung zur Weiterentwicklung der Geistes- und
    Sozialwissenschaften sowie der Gesamtstrategie der Universität Hamburg" des
    Wissenschaftsrats kritisch.

    Der AStA beschäftigt sich mit dem Umzug des Philturms in den Übersee-Ring,
    es wird ein Kriterienkatalog für den Umzug erarbeitet.

    Nächste Woche soll es ein "Stilles Wörtchen" mit dem Schwerpunkt
    Wahlauswertung geben.

    Es gab eine Veranstaltung mit Kerim Schamberger zur Lage in der Türkei nach
    dem Putsch.

    "Uni Hilft" geht in die vierte Runde, jedes Semester können 180 Studierende
    aus den Flüchtlingsunterkünften im Uni-Betrieb integriert werden, der Bedarf
    sei aber deutlich höher.

    Die Veranstaltung mit Farid Esack zur Frage, ob man die Systeme in
    (ehemaligem) Südafrika und Israel vergleichen und als Apartheid kennzeichnen
    könne, sei analytisch und populär gewesen. Sie wird später im Internet
    abrufbar sein.

    Engagement für Verbesserungen im Studienkolleg: die Sprachkurse müssten
    beibehalten werden, die Unterfinanzierung beendet, der Rassismus im
    Unterricht muss aufhören.

    Der Ausschuss gegen Rechts berichtet, dass es am 30. Januar 18.30 Uhr in der
    HWP eine szenische Lesung zur Machtübertragung an die Faschisten, auch
    bezüglich der Kolleg*innen der Universität, geben wird.


    \subsubsection{Anfragen an den AStA}

    Es wird gefragt, ob es keinen Bericht vom Kulturreferat gibt?

    Es wird geantwortet, es gebe intensive Diskussionen darüber, wie das
    Kulturkursprogramm besser organisiert werden kann.

    Es wird kontrovers über die Unterstützung der Veranstaltung mit Farid Esack
    durch das RIS diskutiert. Einige äußern Kritik an seine Person und die
    BDS-Kampagne (Boykott, Disinvestitionen und Sanktionen für Palästina). Andere
    argumentieren, Esack vertrete eine egalitäre Position und die gewaltfreie
    BDS-Kampagne setze sich für eine Friedenslösung im Nahost ein. Einige
    befürchten, dass mit der BDS-Kampagne Stimmung gegen israelische Studierende
    gemacht werden könnte, dagegen wird u.a. für die Verantwortung für Aufklärung
    und ein Ende aller Rüstungsexporte hier gesprochen.

    Das RIS bietet an, den Film von der Veranstaltung gemeinsam zu sehen und
    anschließend darüber zu diskutieren.

    \subsubsection{Dringlichkeitsanträge des AStA}

    Keine.

    \subsubsection{Aktuelle Stunde}

    Golnar beantragt eine aktuelle Stunde zu Trump: Seine Amtszeit sollte
    möglichst nicht sehr lang werden. Die "Radikalen Linken" in den USA haben
    einen Kongreß gemacht: Bürgerrechte sollen stärker mit der sozialen Frage
    verbunden werden.

    In der Diskussion herrscht weitgehend Einigkeit: Gegen die
    milliardärenschwerste Regierung aller Zeiten und die Verbreitung von
    "alternative facts" seien Wissenschaft und Journalismus überall gefordert,
    sich couragiert für die Wahrheit einzusetzen. Es werden Vorschläge zur
    praktischen Solidarität mit US-amerikanischen Wissenschaftler*innen gemacht
    und dafür, wie auch hier die Bedinungen für das Wissenschaftssystem
    verbessert werden können und müssen.

    \newpage
    \subsubsection{Feststellung der endgültigen Fassung des Teils B der Tagesordnung}

    %%% BESCHLOSSENE TAGESORDNUNG %%%
    \underline{Endgültige Fassung der Tagesordnung}

    Franziska möchte die Wahlauswertung vor die englischsprachige Wahlwerbung
    ziehen und den TOP zu VG Wort nach dem Gutachten des Wissenschaftsrats
    behandeln.

    \textit{Campus Grün Fraktionspause bis 20.27 Uhr}

    Die so geänderte Tagesordnung wird mehrheitlich angenommen.

    \begin{enumerate}[label={\textbf{Top \theenumi}},leftmargin=*]
        \item RBCS-Bestätigung (10 Min.)
        \begin{enumerate}
            \item Bestätigung der Wahlniederschrift
            \item Bestätigung der Referent*innen
        \end{enumerate}
        \item Wahlauswertung (60 Min.)
        \item Englischsprachige Wahlwerbung (V1617-048) (30 Min.)
        \item Gutachten des Wissenschaftsrats (V1617-050) (30 Min.)
        \item VG-Wort (V1617-047) (30 Min.)
        \item Verschiedenes (5 Min.)
    \end{enumerate}

    \subsubsection{Feststellung der Beschlussfähigkeit}

    \textit{Mit 25 Parlamentarier*innen beschlussfähig.}

    \subsubsection{Genehmigung der Protokolle der vorangegangenen Sitzungen}

    Anmerkung von Melf: Er erklärt, dass er für die 12. Sitzung entschuldigt war.

    Die Protokolle der 12. und 13. Sitzung werden mit dieser Änderung genehmigt.


    \vspace{0.5cm}
    {\Large \underline{Teil B}}


    \subsection{RBCS-Bestätigung}

    \subsubsection{Bestätigung der Wahlniederschrift}

    Es war eine kleine Versammlung, die Referent*innen wurden wie vorgeschlagen
    gewählt. Geoffrey als Wahlleitung liest die Wahlniederschrift vor.

    Sie wird vom Parlament einstimmig genehmigt.

    \subsection{Bestätigung der Referent*innen}

    Lars und Andrea stellen die Referent*innen Tim Manzella und Dörte Sievert vor,
    die von der RBCS-Vollversammlung gewählt wurden.

    Es gibt eine kurze Diskussion darüber, warum so wenige zur VV gekommen sind,
    obwohl sie sehr gründlich beworben worden sei. Muss die Beratung politischer
    werden? Die Anträge usw. seien sehr aufwendig und die Leute hätten keine Zeit
    für Politik. Es wird sich dafür ausgesprochen, gegen die Ideologie des
    "Klarkommens" und die BaMa-Zumutungen (für alle) die Arbeit für
    Nachteilsausgleich so zu realisieren, daß alle Einschränkungen fallen.

    Die Referent*innen werden einstimmig bestätigt.

    \subsection{Wahlauswertung}

    Aus dem Präsidium: Die Zahlen sind schon vorgetragen, es werden im Nachklapp
    etwa 80 Briefe von der Post abgeholt. Das Präsidium wird eine Auszählung und
    Auswertung der vielen ungültigen Stimmen vornehmen, um herauszufinden, woran
    es liegen kann. Ansonsten seien die Wahlen gut gelaufen, die
    Öffentlichkeitsarbeit hätte besser sein können. Es habe eine Kultur der
    Ernsthaftigkeit gegeben. Fraglich sei, warum am Ende der Wahlauszählung nur
    wenige sich beteiligt haben. Es sei insgesamt viel Arbeit gewesen, dem
    nächsten Präsidium wird empfohlen, rechtzeitig einen Plan zu machen und
    abzuarbeiten.

    Es werden Vorschläge gemacht, wie man die Briefwahl nachvollziehbarer
    gestalten kann, z.B. durch Pictogramme und Bedienungsanleitung auf Video.

    Zum Thema "fairer Wahlkampf" wird gefragt, was von der Kampagne "Wir lieben
    den AStA" zu halten sei? Eine Antwort: Da sei irgendwas mit Mexikanern, was
    man nicht so ernst nehmen müsse.

    Ansonsten sei der Wahlkampf fair und politisch orientiert gewesen, nicht
    viel Zerstörung. Künftig sei mehr dauerhafte Präsenz auf dem Campus, in den
    OEn usw. wichtig. Z.B. könnten die OE-Tutoren zur StuPa-Sitzung eingeladen
    werden. Es wird problematisiert, dass es z.B. in der Biologie keine
    studentischen OE-Tutor*innen (mehr) gibt, da alle vom Studienbüro
    eingestellt sind. Es wird sich für mehr Öffentlichkeit in den eher
    unpolitischen Außenbereichen ausgesprochen. Zu beachten sei auch, nicht
    Fremdwörter zu benutzen, verständlich zu argumentieren und positive
    Forderungen zu stellen. Dagegen wird zu den "Begrifflichkeiten" argumentiert,
    dauerhafte Entwicklung vollziehe sich durch These, Antithese und Synthese und
    die Widersprüche der Gesellschaft ließen sich nicht mit einfacher Sprache
    aufheben.

    Weitere Fragen: Kann mit dem Studierendenwerk geklärt werden, dass die Flyer
    in den Mensen liegenbleiben? Können die Urnen länger aufbleiben, z.B. bis
    18.30 Uhr?

    \subsection{Englischsprachige Wahlwerbung}

    Jim stellt einen Kompromissvorschlag des Präsidiums vor, in dem Vorbehalte und
    Änderungen aus der Diskussion auf der letzten Sitzung berücksichtigt sind.

    In der Diskussion wird nochmals davor gewarnt, dass fehlende Deutschkenntnisse
    ausgenutzt werden könnten (so hätten im Geomatikum vor einigen Jahren Leute
    für die MIN-Liste unterschrieben, weil ihnen erzählt wurde, dass sie sonst
    abgeschoben werden könnten).

    Zu präzisieren sei, dass es sich um die Übersetzung der Listendarstellung
    und nicht um einen anderen Text gehe.

    Hinterfragt wird, was "zumindest" bedeute und warum nur die kleinen
    Listendarstellungen ins Englische übersetzt werden sollen?

    Es wird sich schließlich darauf geeinigt, dass es hier um Mindeststandards
    geht, die erweitert werden können. Das Präsidium müsse entscheiden, was
    darüber hinaus machbar sei.

    Es wird dafür plädiert, auch die allgemeinen Wahl-Informationen und die
    Eidesstattliche Erklärung ins Englisch zu übersetzen.

    Till beantragt, ein weiteres "zumindest" vor "ihre kleinen
    Listendarstellungen" einzufügen.

    Das Präsidium übernimmt den Änderungsantrag.

    Der Antrag wird so einstimmig angenommen.

    \subsection{Gutachten des Wissenschaftsrats ("Empfehlung zur
    Weiterentwicklung der Geistes- und Sozialwissenschaften sowie der
    Gesamtstrategie der Universität Hamburg")}

    Franziska stellt den Antrag des AStA vor: Der Wissenschaftsrat wolle von den
    gesellschaftlichen Aufgaben und der Unterfinanzierung nichts wissen, es gehe
    ihm nur um „Sichtbarkeit“. Die kleine Fächer hätten ihrer Meinung nach
    dagegen hohe Relevanz, z.B. Äthiopisch für die Inklusion von Geflüchteten.

    In der Debatte wird besonders die marktkonforme Orientierung des Gutachtens
    kritisiert und den Wert der geisteswissenschaftlichen Fächer für die Lösung
    der globalen Probleme hervorgehoben.

    In Zeile 14 wird ergänzt: "markt-"konform.

    Der Antrag wird mit m:4:0 angenommen.

    \subsection{VG-Wort}

    Geoffrey stellt den Stand der Überarbeitung des Antrags von ihm und Till vor.  Till erläutert einzelne Aspekte der Überarbeitung und weist besonders auf die im letzten Absatz gefaßten „Dreieinigkeit“ von guter wissenschaftlicher künstlerische Arbeit, dem allgemeine Zugang dazu sowie der sozialen Absicherung der Kultur- und Wissenschaftsschaffenden hin.

    Nach der Klärung einzelner Fragen wir der Antrag vorbehaltlich redaktioneller
    Änderungen mit m:0:6 angenommen.

    \subsection{Verschiedenes}

    \textit{Die Sitzung wird um 22.50 Uhr geschlossen.}
\end{document}
